\clearpage
\section*{Introducere}
\addcontentsline{toc}{section}{Introducere}

Crearea jocurilor este considerată atât o știință cât și o artă reprezintă
procesul de creare a jocurilor video. Dezvoltarea jocurilor video implica
generarea, proiectarea, construirea, testarea și lansarea conceptului. În
timpul creării unui joc, este important de considerat mecanica jocului,
recompensele, implicarea jucătorilor și designul nivelului. Este întreprinsă de
o echipă de dezvoltatori, care poate varia de la o persoană la o mai multe,
acestae fiind angajați în cadrul unei companii sau contribuitori individuali
care dezvoltă jocuri în timpul liber ca hobby. Un membru al acestei echipe ar
putea fi un programator, un designer de sunet, un artist, un designer sau multe
alte roluri disponibile în industrie, cu cât echipa este mai mică cu atât mai
multe roluri va trebui să își asume o singruă persoană.

Jocurile comerciale tradiționale sunt finanțate în mod normal de un editor și
durează câțiva ani pentru a fi dezvoltate. Ca excepție există și jocuri open
source însă acestea sunt făcute de echipe mici, de regulă de o singură
persoană, ele sunt dezvoltate mult mai lent și nu se bucură de o poularitate
mare spre deosebire de jocurile comerciale. Totuși aceste jocuri de obicei nu
sunt făcute cu scopul de a obține un profit, ele sunt făcute din plăcere
sau pentrua învăța ceva nou.

Jocurile video sunt o industrie de miliarde de dolari ceea ce nu este de
mirat considerând faptul că aproximativ trei miliarde de oameni joacă jocuri
video \cite{GamersStatistic}, ceea ce oferă posibilitatea dezvoltatorilor, dar
în special designerilor să influențeze o audiență
vastă. 
Principalele motive ale popularității jocurilor sunt: experiențe unice imposibile în viața reală, combaterea plictiselei și petrecerea timpului cu prietenii la distanță.

Sunt interesat să fac jocuri care utilizează diverși algoritmi și formule
matematice care pot genera conținutul jocului începând de la teren, plante și
terminând cu obiecte și creaturi vii. Totuși aceste lucruri sunt prea
ambițioase pentru o perioadă atât de scurtă și pentru o echipă formată ditr-o
singură persoană, astfel mă voi limita la ceva mai simplu și realizabil în
decursul a 8 săptămâni.
Consultand acest document am ajuns la concluzia ca studiile la UTM mi-au
format o baza buna pentru crearea unor jocuri simple. \cite{GameDevRoadmap}



Pentru dezvoltarea jocurilor sunt folosite mai multe limbaje de programare,
jocurile video fiind dezvoltate din 1962 \cite{FirstGame}, în această perioadă
ele fiind scrise in \textbf{Asembler}. Dezvoltarea jocurilor cu grafica de azi
practic este imposibilă în astfel de limbaje de nivel jos. 

În ziua de azi jocurile s-au dezovltat foarte tare s-au evidențiat o mulțime 
de categorii dintre care următoarele sunt cele mai importante, dar pot exista 
și combinații:
\begin{itemize}
    \item sandbox;
    \item strategie în timp Real (RTS);
    \item shootere (FPS și TPS);
    \item arena de luptă online Multiplayer (MOBA);
    \item jocuri de rol (RPG, ARPG și multe altele);
    \item simulare și sport;
    \item puzzle-uri și jocuri de petrecere;
    \item acțiune-aventură;
    \item supraviețuire și groază;
    \item platforma.
\end{itemize}
